%% start of file `template.tex'.
%% Copyright 2006-2013 Xavier Danaux (xdanaux@gmail.com).
%
% This work may be distributed and/or modified under the
% conditions of the LaTeX Project Public License version 1.3c,
% available at http://www.latex-project.org/lppl/.


\documentclass[12pt,a4paper,sans]{moderncv}        % possible options include font size ('10pt', '11pt' and '12pt'), paper size ('a4paper', 'letterpaper', 'a5paper', 'legalpaper', 'executivepaper' and 'landscape') and font family ('sans' and 'roman')

% moderncv themes
\moderncvstyle{classic}                            % style options are 'casual' (default), 'classic', 'oldstyle' and 'banking'
\moderncvcolor{grey}                              % color options 'blue' (default), 'orange', 'green', 'red', 'purple', 'grey' and 'black'
%\renewcommand{\familydefault}{\sfdefault}         % to set the default font; use '\sfdefault' for the default sans serif font, '\rmdefault' for the default roman one, or any tex font name
%\nopagenumbers{}                                  % uncomment to suppress automatic page numbering for CVs longer than one page

% character encoding
\usepackage[utf8]{inputenc}                       % if you are not using xelatex ou lualatex, replace by the encoding you are using
%\usepackage{CJKutf8}                              % if you need to use CJK to typeset your resume in Chinese, Japanese or Korean

% adjust the page margins
\usepackage[scale=0.75, centering]{geometry}
%\setlength{\hintscolumnwidth}{3cm}                % if you want to change the width of the column with the dates
%\setlength{\makecvtitlenamewidth}{10cm}           % for the 'classic' style, if you want to force the width allocated to your name and avoid line breaks. be careful though, the length is normally calculated to avoid any overlap with your personal info; use this at your own typographical risks...
\usepackage{ragged2e}
% personal data
\name{Nicola}{Vianello}
%\title{Resumé title}                               % optional, remove / comment the line if not wanted
\address{Corso Stati Uniti 4}{I-35127 Padova}{Italy}% optional, remove / comment the line if not wanted; the "postcode city" and and "country" arguments can be omitted or provided empty
%\phone[mobile]{+1~(234)~567~890}                   % optional, remove / comment the line if not wanted
\phone[fixed]{+39~(049)~829~5991}                    % optional, remove / comment the line if not wanted
%\phone[fax]{+3~(456)~789~012}                      % optional, remove / comment the line if not wanted
\email{nicola.vianello@igi.cnr.it}                               % optional, remove / comment the line if not wanted
% \homepage{www.johndoe.com}                         % optional, remove / comment the line if not wanted
% \extrainfo{additional information}                 % optional, remove / comment the line if not wanted
% \photo[64pt][0.4pt]{picture}                       % optional, remove / comment the line if not wanted; '64pt' is the height the pictu
% re must be resized to, 0.4pt is the thickness of the frame around it (put it to 0pt for no frame) and 'picture' is the name of the picture file
% \quote{Some quote}                                 % optional, remove / comment the line if not wanted

% to show numerical labels in the bibliography (default is to show no labels); only useful if you make citations in your resume
%\makeatletter
%\renewcommand*{\bibliographyitemlabel}{\@biblabel{\arabic{enumiv}}}
%\makeatother
%\renewcommand*{\bibliographyitemlabel}{[\arabic{enumiv}]}% CONSIDER REPLACING THE ABOVE BY THIS

% bibliography with mutiple entries
%\usepackage{multibib}
%\newcites{book,misc}{{Books},{Others}}
%----------------------------------------------------------------------------------
%            content
%----------------------------------------------------------------------------------
\begin{document}
%-----       letter       ---------------------------------------------------------
% recipient data
\recipient{EUROfusion Program Management
  Unit}{EUROfusion\\Boltmanstra{\"s}{\"s}e 2 Garching,  Germany}
\date{June 12, 2020}
\opening{Dear Sir/Madam,}
\closing{Yours faithfully, \\
  \vspace{2pt}
  \includegraphics[height=4\baselineskip]{signature.pdf}}
%\enclosure[Attached]{curriculum vit\ae{}}          % use an optional argument to use a string other than "Enclosure", or redefine \enclname
\makelettertitle
\justify
my name is Nicola Vianello, I am a 44 years old Phd Physic Scientist,
currently working at the Consorzio RFX, Padova. I would like to apply for
the position of Task Force Leader for Tokamak Exploitation within the Fusion Science
Department. 

I've been involved in Fusion plasma Science since my M.Sci. Thesis in
Physics in 1999. My primary research interest is transport phenomena in fusion
oriented plasmas with strong emphasis on non-linear dynamics. I have
approached the problem both experimentally, through the collection, 
analysis, interpretation and modeling of experimental data, and
numerically, through the use of massive parallel fluid codes. 

I've a strong attitude in
data analysis and interpretation with a particular emphasis on the
comparison with theories and codes used as a suitable framework for the
correct interpretation of real data. I've been involved in
the investigation of a variety of phenomena, ranging from electrostatic
turbulence particle transport, sheared flow turbulence generation,
turbulent non-linearly generated structures, 3D
magnetic perturbation and its relation with plasma transport phenomena, ELMs, divertor
physics and the connection between divertor and upstream SOL and
pedestal properties. I've worked on a variety of different
devices with the rather unique opportunity to extend my competences
on all the major magnetic configurations, tokamaks,  stellarators and
reversed field pinches. Addressing each of theses scientific 
topics has required the development of a solid theoretical
background, and the adaptation of theories to the framework
studied. Motivated by my scientific curiosity I've established
fruitful collaborations with the plasma astrophysics community: this
experience helped me to keep a broad-minded approach for a deeper mastery of the basic 
phenomena shared by the two plasma environments as turbulence or particle energization. 
I strongly believe that my scientific carrier testifies my 
attitude towards a comprehensive integrated understanding of plasma and tokamak physics. 

I've been involved in the WP-MST1 as Scientific Coordinator since
its establishment (AUG14-2.2-3, TCV15-2.2-3, TCV15-1.5-1, Topic 21 and
Topic 16) but I have been as well an active contributor in WP-JET1, acting as Scientific Coordinator
of Task T18-02 and Experiment M18-41. The developed expertise aligns
to those required for Mission 2 of the EUROfusion Roadmap, for the
determination of a viable plasma exhaust solution, as well as
for Mission 1 given the interest on the interplay between divertor
conditions and upstream SOL and pedestal. 
I've enjoyed the
possibility offered by the EUROfusion framework to work on
different devices. I strongly believe that the level of
comprehension mandatory to extrapolate to future
machines requires an appropriate comparative approach, overcoming 
specific issues of singular experiments.
EUROfusion offered me the possibility to work within a stimulating
scientific environment. I believe that gathering
european expertise to work on specific issues in a
goal-oriented approach helped Europe as a whole community to
consolidate the leadership on fusion science.

The role of Scientific Coordinator induce me to increase my  management and coordination capabilities in an
international environment. I've always acted in an inclusive
manner, by properly spreading knowledge and information while keeping
the leadership in pushing forward the scientific objectives. I'm a
strong supporter of Open Science and I've promoted it
within the scientific team in the hope of speeding up the research activity.

The role of TFL within the Tokamak Exploitation workpackage would represent
for me a stimulating and aspiring challenge. The TE workpackage, in
particular for the Scientific program concerning the Medium Size
Tokamaks, will focus on extending even further our present understanding of
the tokamak physics providing robust scenarios in term of performance
and exhaust solution for DEMO class devices, while consolidating
the ITER physics understanding to increase the confidence for a
safe and successful operation of ITER. As aforementioned this is an
ambitious task which require a focused program while keeping an open
minded approach to extend our knowledge towards partially unknowns
territories. I think my competences fits well with the requirements
for the position and I hope I could be able to provide my personal contribution.


\makeletterclosing

\end{document}


%% end of file `template.tex'.



% \documentclass[12pt,stdletter,a4paper,dateno,sigleft]{newlfm}
% % \usepackage{kpfonts}
% \usepackage{url}
% %\usepackage{charter}
% \usepackage{fontspec,lipsum}
% % ----------
% % font used
% \usepackage{fontspec,lipsum, xltxtra, xunicode}
% \defaultfontfeatures{Mapping=tex-text}
% \newfontfamily\bodyfont[]{Helvetica Neue}
% \newfontfamily\thinfont[]{Helvetica Neue UltraLight}
% \newfontfamily\headingfont[]{Helvetica Neue Condensed Bold}
% \setmainfont[Mapping=tex-text, Color=textcolor]{Helvetica Neue Light} 

% %\widowpenalty=1000
% %\clubpenalty=1000

% %for the logo
% \newsavebox{\RFXlogo}
% \sbox{\RFXlogo}{%
% 	\parbox[b]{1.75in}{%
% 		\vspace{0.5in}%
% 		\includegraphics[scale=.2,ext=.pdf]
% 		{RFX-logo}%
% 	}%
% }%
% \makeletterhead{Uiuc}{\Lheader{\usebox{\RFXlogo}}}

% % for the signature
% \newsavebox{\Sigx} 
% \sbox{\Sigx}{%
%           \includegraphics[height=7\baselineskip]{signature.pdf}}
% \signature{\usebox{\Sigx}}
% \makesignature{NV}{\newsavebox{\Signature}} 

% \newlfmP{headermarginskip=10pt}
% \newlfmP{sigsize=50pt}
% \newlfmP{dateskipafter=20pt}
% %\newlfmP{paperheight=700pt}
% \newlfmP{addrfromphone}
% \newlfmP{addrfromemail}
% \newlfmP{MinFoot=13pt}
% %\newlfmP{MinHead=50pt}
% \PhrPhone{Phone}
% \PhrEmail{Email}
% %\encllist{Reference Letter from Prof. Francesco Gnesotto}
% % \lthUiuc
% \sigNV

% \namefrom{Nicola Vianello}
% \addrfrom{%
%     \today\\[10pt]
%     Consorzio RFX \\
%     Corso Stati Uniti 4\\
%     I-35127 Padova\\
%     Italy
% }
% \phonefrom{+39 0498295991}
% \emailfrom{nicola.vianello@igi.cnr.it}

% %\addrto{%
% %\textbf{To:} Mr. Hans Jahreiss, Fusion for Energy, European Domestic Agency \\
% %\textbf{Job Title:} Plasma Physics,  Scientific Coordinator \\
% %\textbf{Ref:}  Senior Scientific Officer POP-009 }

% \greetto{To Whom It May Concern,}
% \closeline{Sincerely,}
% \begin{document}
% \begin{newlfm}
% my name is Nicola Vianello, I am a 40 years old Phd Physic Scientist,
% currently working at the Consorzio RFX, Padova. I would like to apply for
% the ITPA position in the Pedestal and Edge physics group. 

% I've been involved in Fusion plasma Science since my M.Sci. Thesis in
% Physics in 1999. My primary research interests is transport phenomena in fusion
% oriented plasmas with strong emphasis on non-linear dynamics. I have
% addressed the problem both experimentally, through the collection, 
% analysis, interpretation and modeling of experimental data, and
% numerically, through the use of massive parallel fluid codes. 

% I've a strong attitude in
% data analysis and evaluation with a particular emphasis on the
% comparison with theories and codes which provide the suitable framework for the
% correct interpretation of real data. I've been involved in
% the interpretation of a variety of phenomena, ranging from electrostatic
% turbulence particle transport, sheared flow turbulence generation,
% turbulent non-linearly generated structures, interplay between 3D
% magnetic perturbation and plasma transport phenomena, relationship
% between divertor conditions and upstream SOL and pedestal properties. Each of these
% topics has required the development of a solid theoretical
% background, and the adaptation of theories to the studied framework. I strongly believe that this
% attitude, which combines experimental expertise  with clear and solid
% theoretical background provide an additional important element for the
% establishment of a comprehensive integrated comprehension of
% tokamak and more generally plasma physics. I strongly believe that a
% physical properties of the pedestal and edge confined region are
% tightly linked to the SOL and boundary condition: moving from this
% assumption I'm trying to focus my experimental activities to bridge
% the gap between these two region and I believe that this expertise may
% be extremely relevant for the Pedestal and Edge physics group. 
% I've been actively involved in the coordinations of experimental
% activities in various international experiment (RFX-mod, JET, TCV and
% Asdex-Upgrade), both as task force leader and as scientific
% coordinator. I believe these experiences
% provide me a good management and coordination capabilities also in an
% international environment. My international experience is indeed very good, with vital and active
% collaborations with different European and international laboratories. 

% I think that during my research carrier I have proved strong autonomy
% accompanied by good capability to work in small and large groups. 

% All these qualities and competences fits well with the
% requirements for the position. 

% \end{newlfm}

% \end{document}
